% Abstrak (Indonesia dan Inggris): satu paragraf, maksimal 300 kata, satu spasi. Ganti teks dalam [ ... ] dengan isi penelitian Anda.
\babtp{ABSTRAK}
\begin{abstrak}
Penelitian ini dilatarbelakangi oleh [permasalahan utama yang dihadapi dan alasan penelitian perlu dilakukan]. Berdasarkan hasil observasi awal dan kajian pustaka, [uraikan kesenjangan antara kondisi yang diharapkan dan kondisi nyata, disertai data atau sitasi pendukung]. Hasil analisis awal menunjukkan bahwa [temuan awal atau studi pendahuluan], sehingga penelitian ini dipandang layak dan mendesak untuk dilaksanakan. Apabila tidak ditangani, [dampak yang mungkin timbul] akan semakin meluas, sehingga diperlukan penelitian yang [karakteristik penelitian yang dibutuhkan]. Permasalahan yang dirumuskan dalam penelitian ini adalah [rumusan masalah pertama], [rumusan masalah kedua], dan [rumusan masalah ketiga]. Tujuan penelitian ini adalah [tujuan penelitian secara ringkas, sesuai rumusan masalah]. Penelitian dilaksanakan di [lokasi penelitian] selama [lama penelitian] dengan menggunakan metode [nama metode atau pendekatan]. Tahapan penelitian meliputi [tahap pertama], [tahap kedua], dan [tahap ketiga]. Data dikumpulkan melalui [teknik pengumpulan data] dan dianalisis dengan [teknik analisis data]. Perancangan sistem dilakukan dengan membagi sistem menjadi [jumlah] bagian utama, yaitu [bagian pertama] yang berfungsi untuk [fungsi bagian pertama], [bagian kedua] yang berfungsi untuk [fungsi bagian kedua], serta [bagian ketiga] yang berfungsi untuk [fungsi bagian ketiga]. Sistem atau objek yang diteliti dirancang menggunakan [perangkat keras dan perangkat lunak yang digunakan]. Pengujian dilakukan dengan [jumlah] skenario, yaitu [skenario pertama], [skenario kedua], dan [skenario ketiga], masing-masing diulang sebanyak [jumlah pengulangan] kali agar hasilnya dapat diandalkan. Kinerja diukur berdasarkan [indikator kinerja], dengan parameter yang diamati meliputi [parameter pertama], [parameter kedua], dan [parameter ketiga] serta nilai acuan [nilai acuan atau standar yang digunakan]. Hasil penelitian menunjukkan bahwa [hasil utama pertama beserta nilainya] dan [hasil utama kedua beserta nilainya]. Berdasarkan hasil tersebut disimpulkan bahwa [kesimpulan utama yang menjawab rumusan masalah]. Keterbatasan penelitian ini adalah [batasan penelitian], sehingga hasilnya terutama berlaku pada [cakupan hasil penelitian]. Penelitian ini diharapkan bermanfaat bagi [pihak yang memperoleh manfaat], dan disarankan [saran pengembangan penelitian selanjutnya].

\textbf{Kata kunci:} kata kunci 1, kata kunci 2, kata kunci 3.
\end{abstrak}

\begin{halamanbaru}\end{halamanbaru}
\babtp{ABSTRACT}
\begin{abstrak}
This research is motivated by [the main problem and the reason the study is needed]. Based on preliminary observations and a literature review, [describe the gap between the expected and actual conditions, supported by data or citations]. A preliminary analysis shows that [the initial finding or preliminary study], so this research is considered feasible and urgent. If left unaddressed, [the possible impact] will continue to spread, so research that [the required characteristics of the study] is needed. The problems formulated in this research are [first problem], [second problem], and [third problem]. The objective of this research is [the research objective, consistent with the problem statement]. The research was carried out at [research location] for [duration] using [method or approach]. The stages of the research include [first stage], [second stage], and [third stage]. Data were collected through [data collection technique] and analyzed using [data analysis technique]. The system under study was designed using [hardware and software]. Testing was conducted in [number] scenarios, namely [first scenario], [second scenario], and [third scenario], each repeated [number of repetitions] times to ensure reliable results. Performance was measured based on [performance indicators], with the observed parameters including [first parameter], [second parameter], and [third parameter] and the reference value of [reference value or standard used]. The results show that [first main result and its value] and [second main result and its value]. Based on these results, it is concluded that [main conclusion answering the problem statement]. The limitation of this research is [research limitation], so the results mainly apply to [scope of the results]. This research is expected to benefit [beneficiaries], and further research is suggested on [suggestions for further research].

\textbf{Keywords:} keyword 1, keyword 2, keyword 3.
\end{abstrak}

\begin{halamanbaru}\end{halamanbaru}
